\section{The complete reflected representation}\label{sec:representation}

\subsection{An exact identity before asymptotics}

Let
\[
 s_\pm(z)=\frac{1\pm\ii z}{2},\qquad
 M(s)=\frac{s(s-1)}{16}\pi^{-s/2}\Gamma(s/2),\qquad
 U^*(s)=\overline{U(\bar s)}.
\]
The exact Riemann--Siegel contour remainder is
\begin{equation}\label{eq:contour}
 R_{0,N}(s)=M(s)\int_{N\swarrow N+1}
       \frac{w^{-s}\ee^{\ii\pi w^2}}{\ee^{\pi\ii w}-\ee^{-\pi\ii w}}\dd w.
\end{equation}
The contour crosses \((N,N+1)\) and has the orientation specified in
\cite[Section 4, equations (36)--(38)]{polymath}. For \(T>0\) and
nonreal \(s\), put
\begin{align}
 r_{T,n}(s)&=\frac1{\sqrt\pi}\int_\R
  \ee^{-v^2}M(s+\sqrt T\,v)n^{-s-\sqrt T\,v}\dd v,\\
 R_{T,N}(s)&=\frac1{\sqrt\pi}\int_\R
  \ee^{-v^2}R_{0,N}(s+\sqrt T\,v)\dd v,\qquad
 U_{T,N}(s)=\sum_{n\le N}r_{T,n}(s)+R_{T,N}(s).
\end{align}
For every fixed integer \(N\ge0\) and \(\Rea z\ne0\),
\begin{equation}\label{eq:exactRS}
 H_T(z)=U_{T,N}(s_+(z))+U_{T,N}^*(s_-(z)).
\end{equation}
To see the heat step, insert the Gaussian identity
\[
 \ee^{Tu^2}\cos(zu)=\frac1{\sqrt\pi}\int_\R
       \ee^{-v^2}\cos((z-2\ii\sqrt T\,v)u)\dd v
\]
in \eqref{eq:flow}, and then insert the exact two-contour formula for
\(H_0\). The Gaussian dominates the sub-Gaussian growth of each
complete arm on the permitted shifted lines. The same local domination
justifies differentiation:
\begin{equation}\label{eq:exactRSprime}
 H_T'(z)=\frac\ii2U_{T,N}'(s_+(z))
              -\frac\ii2(U_{T,N}^*)'(s_-(z)).
\end{equation}
Both the gamma factor and the full contour remainder are differentiated.
Off the real axis, \eqref{eq:exactRS} is not twice a real part.

On the real axis, let \(\mathscr A(x)=U_{T,N}(s_+(x))\).
Then \(H_T=2\Rea\mathscr A\). If \(\mathscr A=\rho\ee^{\ii\psi}\ne0\)
at a zero, differentiation gives
\[
 H_T'=-2\rho\psi'\sin\psi.
\]
A multiple zero therefore requires \(\psi'=0\), or belongs to the
separate alternative \(\mathscr A=0\). No strict phase rotation or
nonzero arm floor follows from the reflected representation alone.
Indeed \(\mathscr A\mapsto\mathscr A+\ii q\), with \(q\) a real-entire
heat solution, preserves the reduced data \(H=\mathscr A+\mathscr A^*\).
This is an information test, not permission to alter the canonical arm.
Section~\ref{sec:controls} gives an exact threshold control.

\subsection{The named approximation input and its domain}

The asymptotic work uses the internal complex-strip theorem HDA CS3/CS6
and its application GNP-RS E01 \cite{hda,gnpe01}.
We state exactly the portion consumed here. It is not re-proved as an
independent provider theorem in this manuscript.

On a fixed chart about a large positive real center \(a\), choose the
branches continued from \(a\), and hold
\(N=\lfloor\sqrt{a/(4\pi)}\rfloor\) fixed. Define
\begin{align}
 L(z)&=\Log(z/(4\pi)),\\
 A(T,z)&=\frac{\pi^{1/4}}{16}z^{7/4}
   \exp\left[-\frac{\pi z}{8}
       +\frac T{16}\left(L(z)^2-\frac{\pi^2}{4}\right)\right],
       \label{eq:carrier}\\
 \Theta(T,z)&=\frac z4(L(z)-1)+\frac{\pi T L(z)}{16}
                         +\frac{7\pi}{8},\qquad
 \nu(z)=\frac z2+\frac{\pi T}{8},\\
 C_\pm(T,z)&=\frac{A(T,z)}2\ee^{\pm\ii\Theta(T,z)},\\
 w_n(T,z)&=n^{-1/2}
  \exp\left\{\frac T4\bigl((\log n)^2-L(z)\log n\bigr)\right\}.
\end{align}
The retained correction is
\begin{align}
 \rho(z)&=\sqrt{z/(4\pi)},\qquad
 G(z)=\bigl[M(s_+(z))M(s_-(z))\bigr]^{1/2},\\
 \mathcal C(p)&=\frac{\cos(2\pi(p^2-p)-\pi/8)}{\cos(2\pi p)},\qquad
 \mathcal R_N(z)=(-1)^N G(z)\rho(z)^{-1/2}\mathcal C(\rho(z)-N).
 \label{eq:Rcorrection}
\end{align}
The branch of \(G\) is positive on the real segment. Apparent poles of
\(\mathcal C\) are removable. This is the exact-gamma correction in the
approximant; it must not be confused with the full exact contour
remainder \(R_{T,N}\) in \eqref{eq:exactRS}.
Put
\begin{equation}\label{eq:FN}
 F_N(T,z)=A(T,z)\sum_{n\le N}w_n(T,z)
       \cos\bigl(\Theta(T,z)-\nu(z)\log n\bigr)
       +\ee^{-\pi^2T/64}\mathcal R_N(z).
\end{equation}

\begin{inputthm}[Buffered complete field estimate]\label{input:provider}
For each declared fixed target rectangle, there are finite constants and
an onset such that, for \(0\le T\log(a/(4\pi))\le8\),
\begin{equation}\label{eq:provider}
 |H_T(z)-F_N(T,z)|\le C a^{-3/4}(|C_+(T,z)|+|C_-(T,z)|).
\end{equation}
The difference is holomorphic on the chart. The input is uniform in
fractional-cutoff position. It includes the two complete populations,
the retained correction, the enlarged source charts, and the entire
Gaussian exterior.
\end{inputthm}

For the present application one may take target horizontal radius \(1\)
and vertical height \(9\), retaining the larger source buffers of
\cite{gnpe01}. Its tail constants \(c_9,P_9\) are fixed before choosing
\[
 V_9\ge\max\{16,17,16c_9,8\sqrt{P_9+6}\}.
\]
The required initial charts have radius \(V_9+3\), height \(10\), and
radius \(2\), height \(10+2V_9\), with their distinct onsets.
This records the dependence that permits Cauchy estimates later.
No unbounded Gaussian integral is estimated from a bound valid only on
a narrow strip.

Polymath's effective value Theorem 1.3 has its own printed domain
\(x\ge200,\ 0\le y\le1,\ 0<T\le1/2\). We do not enlarge that domain
or differentiate that theorem to obtain our buffered estimates.
The broader complex-strip statement is the named internal input above,
at its existing evidence level.
